% =========================================================================
\section{Introduction}
\label{sec:intro}

Whether a public blockchain is ``decentralized'' is routinely settled by
citing a single concentration statistic: a Gini coefficient of token
holdings, a top-$k$ share, a Nakamoto coefficient, or a
Herfindahl--Hirschman Index (HHI). Early whitepapers
\cite{nakamoto2008bitcoin, wood2014ethereum} described public ledgers as
engines of disintermediation, and an empirical literature has since asked
how far practice has departed from that description
\cite{srinivasan2017quantifying, sai2021taxonomy, ovezik2022sok,
campajola2022evolution}. Token ledgers in particular are reported to be
extremely unequal: the largest ERC-20 tokens have Gini coefficients close
to one \cite{sai2021wealth}, and decentralized-finance governance tokens are
held by a small number of addresses once custodial contracts are mapped to
their beneficiaries \cite{nadler2022defi}.

Each of these numbers is the output of a measurement that requires choices
the ledger does not make for the analyst. Which holders count? Are burn
addresses, issuer reserves and protocol contracts market participants? Is
the unit of analysis an address or the entity that controls it? Over what
horizon are flows aggregated? Which transfers are data and which are spam?
How are repeated transfers between the same pair of addresses represented?
None of these questions has an answer that the data dictate, and every
answer is defensible. This paper asks how much they matter, on two ERC-20
ledgers observed over an identical 90-day window: the Chainlink (LINK)
oracle token and the Uniswap (UNI) governance token.

\subsection{Research questions}
\label{sec:rq-list}
We organise the analysis around three questions, each answered by one
contribution.
\begin{itemize}
  \item \textbf{RQ1 (specification).} On a single ledger, how far do the
  discretionary choices of a concentration measurement---holder population,
  address type, entity resolution, temporal aggregation, flow direction,
  data cleaning and graph construction---move the verdict, and does the
  Gini coefficient move with the HHI?
  \item \textbf{RQ2 (hidden brokers).} A parameterised criterion selects
  addresses that are structurally central, hold zero or dust balances, and
  carry no label in a reference label set. What does the selection depend on: the
  representation used to score centrality, the counterparty count, the
  data-cleaning specification, or the label set?
  \item \textbf{RQ3 (layers).} How are the ownership, routing and rule
  layers of a token ledger related, and does the rule layer have the
  authority to change the routing layer?
\end{itemize}

\subsection{Contributions}

\paragraph{C1: Specification dependence of on-chain concentration}
Three choices act on balance concentration (holder population, address
type, entity resolution), two on flow concentration (temporal aggregation,
flow direction), and two implementation choices act on everything
downstream (data cleaning, edge aggregation). Over the eighteen
combinations of the three balance choices (\ref{app:cleaning}) the
balance HHI of UNI ranges from \V{combo.uni.nf_hhi_min} to
\V{combo.uni.nf_hhi_max}, a factor of \V{combo.uni.nf_hhi_ratio}, and of
LINK from \V{combo.link.nf_hhi_min} to \V{combo.link.nf_hhi_max}, while the
balance Gini moves by at most \V{combo.uni.gini_range_nofloor} and
\V{combo.link.gini_range_nofloor}. The choices act on different ends of the
distribution. Address type and entity resolution move the top: on either
ledger they change the HHI by up to a factor of \V{fx.top.hhi_max} and the
Gini by at most \V{fx.top.gini_max}. A balance floor moves the tail: a USD~100
floor raises $1-\mathrm{Gini}$ \V{link.omg.ratio}-fold on LINK and
\V{uni.omg.ratio}-fold on UNI and changes the HHI by at most
\V{floor.hhi_change_max}. Restricting the population to transfer-active
addresses moves both (Table~\ref{tab:spec}). Two of the rows
are not matters of taste but of implementation. A rule-based removal of a
coordinated dust airdrop on LINK---\V{link.airdrop_tx} transfers from
\V{link.airdrop_senders} senders---removes \V{link.active_drop_pct} of the
addresses that would otherwise count as transfer-active and moves the
balance HHI of that population from \V{link.hhi.active_noair} to
\V{link.hhi.active}. And whether repeated transfers between an address pair
are summed or overwritten, a one-line choice in a standard graph library,
discards \V{link.ew.loss_1w} of transferred volume in the activity-ranked
subgraph and changes which addresses the
broker criterion selects (Section~\ref{sec:edgecase}). We set out
which published studies use which specification
(Table~\ref{tab:lit}) and recommend that on-chain concentration be reported
as a specified range.

\paragraph{C2: A parameterised criterion for hidden brokers, and what it
depends on}
A \emph{hidden broker} is an address that, under a structural score $s$ and
a reference label set $L$, falls in the top percentile of $s$, holds no
more than the median balance, has at least five counterparties and
non-zero sampled betweenness, and is unlabelled in $L$. Stating $s$, $L$
and the data specification as parameters is the point: we show how the
selected set moves with each. Under the main specification the criterion
selects \V{link.brokers} addresses on LINK and \V{uni.brokers} on UNI,
\V{shared.brokers} of them common to both. A counterparty-matched control
drawn from the same eligibility pool and verified blind tests whether the
representation used to compute $s$ adds anything beyond the counterparty
condition (Section~\ref{sec:control}). It does not detect any: matched
controls resolve to named routing protocols about as often as the selected
brokers.

\paragraph{C3: Three layers, one of them able to rewrite the others}
Across addresses, balance and structural position are only weakly
rank-correlated (Section~\ref{sec:wealth-power}), which is how a Gini near
one coexists with an unconcentrated routing layer. On UNI we add a rule
layer measured in delegated votes, the quantity Uniswap governance actually
counts. At the end of the window \V{uni.gov.deleg_share} of the supply was
delegated; \V{uni.gov.n_quorum} delegates sufficed to reach the
$4\times10^{7}$ UNI quorum; and none of the \V{uni.gov.n_eligible}
delegates able to submit a proposal appears among the hundred largest
senders of UNI. The rule layer is not hypothetical: three months before the
window, governance activated a protocol fee and burned $10^{8}$ UNI from
the treasury, and the fee mechanism burned a further \V{uni.burn.window}
UNI during the window (Section~\ref{sec:governance}).

The remainder of the paper is organised as follows.
Section~\ref{sec:related} reviews related work and Section~\ref{sec:data}
describes the data, the cleaning rules and the graph construction.
Sections~\ref{sec:spec}, \ref{sec:brokers} and~\ref{sec:layers} answer
RQ1--RQ3 in turn, Section~\ref{sec:airdrop} documents the airdrop that the
cleaning rule removes, and Sections~\ref{sec:limitations}
and~\ref{sec:conclusion} discuss limitations and conclude.
